\documentclass[12pt,oneside]{book}

% ===================== METADATOS =====================
\newcommand{\universidad}{Universidad de Buenos Aires (UBA)}
\newcommand{\facultad}{Facultad de Derecho}
\newcommand{\materia}{Derechos Reales}
\newcommand{\titulo}{Publicidad Inmobiliaria y Registro de Bienes Inmuebles}
\newcommand{\autor}{Isaac Edgar Camacho Ocampo}
\newcommand{\anio}{2026}

% ===================== PAQUETES =====================
\usepackage[utf8]{inputenc}
\usepackage[T1]{fontenc}
\usepackage[provide=*,spanish]{babel}
\usepackage[a4paper,top=2.5cm,bottom=2.5cm,left=3cm,right=2.5cm,headheight=15pt]{geometry}
\usepackage{amsmath}
\usepackage{amssymb}
\usepackage{mathtools}
\usepackage{graphicx}
\usepackage{float}
\usepackage{caption}
\usepackage{subcaption}
\usepackage{booktabs}
\usepackage{array}
\usepackage{longtable}
\usepackage{tabularx}
\usepackage{enumitem}
\usepackage{xcolor}
\usepackage[most]{tcolorbox}
\usepackage{fancyhdr}
\usepackage{ragged2e}
\usepackage[
    colorlinks=true,
    linkcolor=blue!50!black,
    citecolor=green!40!black,
    urlcolor=blue!60!black,
    pdftitle={\titulo},
    pdfauthor={\autor},
    pdfsubject={Apuntes universitarios},
    bookmarks=true
]{hyperref}

% ===================== CONFIGURACIÓN =====================
\graphicspath{{./img/}{./graficos/}{./figuras/}}
\setcounter{tocdepth}{2}
\setlength{\parindent}{0pt}
\setlength{\parskip}{0.5em}
\setlist[itemize]{topsep=0.3em,itemsep=0.2em}
\setlist[enumerate]{topsep=0.3em,itemsep=0.2em}
\clubpenalty=10000
\widowpenalty=10000

% ===================== COMANDOS =====================
\newcommand{\abs}[1]{\left|#1\right|}
\newcommand{\norm}[1]{\left\lVert#1\right\rVert}
\newcommand{\vect}[1]{\mathbf{#1}}
\newcommand{\deriv}[2]{\frac{d#1}{d#2}}
\newcommand{\pderiv}[2]{\frac{\partial#1}{\partial#2}}

% ===================== CAJAS =====================
\newtcolorbox{definition}[1]{enhanced,breakable,title={Definición: #1},fonttitle=\bfseries,colback=blue!3,colframe=blue!50!black,boxrule=0.8pt,arc=2mm}
\newtcolorbox{important}{enhanced,breakable,title={Importante},fonttitle=\bfseries,colback=yellow!8,colframe=orange!70!black,boxrule=0.8pt,arc=2mm}
\newtcolorbox{example}[1]{enhanced,breakable,title={Ejemplo: #1},fonttitle=\bfseries,colback=green!4,colframe=green!40!black,boxrule=0.8pt,arc=2mm}
\newtcolorbox{formula}[1]{enhanced,breakable,title={Fórmula / Regla: #1},fonttitle=\bfseries,colback=purple!4,colframe=purple!50!black,boxrule=0.8pt,arc=2mm}
\newtcolorbox{exercise}[1]{enhanced,breakable,title={Ejercicio: #1},fonttitle=\bfseries,colback=gray!5,colframe=gray!60!black,boxrule=0.8pt,arc=2mm}
\newtcolorbox{note}{enhanced,breakable,title={Nota},fonttitle=\bfseries,colback=white,colframe=gray!50,boxrule=0.6pt,arc=2mm}

% ===================== CONFIGURACIÓN FINAL =====================
\pagestyle{fancy}
\fancyhf{}
\fancyhead[L]{\nouppercase{\leftmark}}
\fancyhead[R]{\materia}
\fancyfoot[C]{\thepage}
\renewcommand{\headrulewidth}{0.4pt}
\renewcommand{\footrulewidth}{0pt}
\fancypagestyle{plain}{\fancyhf{}\fancyfoot[C]{\thepage}\renewcommand{\headrulewidth}{0pt}}

\begin{document}

% ===================== PORTADA =====================
\begin{titlepage}
\thispagestyle{empty}
\begin{center}
\vspace*{1cm}
{\large \textsc{\universidad}}\\[0.2cm]
{\large \textsc{\facultad}}

\vspace{3cm}
{\Huge \textbf{\materia}}

\vspace{1cm}
{\LARGE \titulo}

\vspace{0.8cm}
\textit{\small Comprender los conceptos es el primer paso para convertir el estudio en conocimiento.}

\vspace{1.2cm}
\rule{0.70\textwidth}{0.5pt}

\vspace{0.5cm}
{\large Apuntes de estudio}

\vspace{0.5cm}
\rule{0.70\textwidth}{0.5pt}

\vfill
{\large \autor}

\vspace{0.3cm}
{\large \anio}
\end{center}
\vfill
\begin{flushleft}
\textit{Este es un modesto aporte a los estudiantes de ``\materia'' no pretende sustituir el material oficial de la materia.}
\end{flushleft}
\end{titlepage}

% ===================== ÍNDICE =====================
\tableofcontents

% ===================== CONTENIDO =====================
\chapter*{Introducción: importancia del tema para el profesional del derecho}
\addcontentsline{toc}{chapter}{Introducción: importancia del tema}
\markboth{Introducción}{Introducción}

Los derechos reales inmobiliarios y su publicidad constituyen pilares fundamentales del tráfico jurídico en Argentina. Su dominio permite al profesional del derecho:

\begin{itemize}
    \item \textbf{Proteger los derechos de los clientes} en transacciones inmobiliarias, evitando fraudes y garantizando los títulos.
    \item \textbf{Comprender el sistema registral argentino (Ley 17.801)}, que regula cómo se adquieren, transmiten y extinguen derechos sobre inmuebles.
    \item \textbf{Prevenir daños} en operaciones inmobiliarias. El caso de la ``sombrilla registral'' (bloqueo registral) muestra cómo una inscripción tardía puede costarle todo al cliente.
    \item \textbf{Aplicar los 8 principios registrales} (rogación, inscripción, presunción, especialidad, legalidad, prioridad, tracto sucesivo, publicidad) en la práctica diaria.
    \item \textbf{Asesorar correctamente} sobre embargos, medidas cautelares y protección de patrimonios.
    \item \textbf{Trabajar eficientemente con escribanos y registros de la propiedad}, comprendiendo los plazos, requisitos y excepciones del sistema.
\end{itemize}

Este conocimiento no es solo teórico: la seguridad del tráfico inmobiliario depende de profesionales que entiendan cómo funciona la publicidad registral. Los clientes necesitarán garantizar que sus derechos reales sean oponibles a terceros, que las inscripciones se realicen dentro de los plazos críticos y que sus patrimonios estén protegidos.

% ---------------------------------------------------------------
\chapter{Derecho real y publicidad inmobiliaria}
\markboth{Derecho real y publicidad}{Derecho real y publicidad}

\section{Introducción al derecho real y la publicidad inmobiliaria}

Hablar de publicidad en materia de derechos reales es hablar propiamente del derecho real. Entre las diferencias esenciales entre derechos reales y derechos personales se encuentra que \textbf{el derecho real es absoluto y oponible \emph{erga omnes}}, mientras que el derecho personal es relativo y resulta oponible únicamente a las partes que integran el vínculo obligacional: el sujeto activo y el sujeto pasivo vinculado por la prestación.

El derecho real es oponible \emph{erga omnes} (a todos) y debe ser respetado por todos. Si debe ser respetado por todos, necesariamente debe ser conocido por todos, porque \textbf{nadie puede ser obligado a respetar aquello que no conoce}. Por ello, la publicidad es una exigencia intrínseca y esencial de la inherencia del derecho real.

\begin{important}
El derecho real es absoluto y oponible a todos; por eso exige publicidad. El derecho personal es relativo y solo es oponible a las partes del vínculo obligacional.
\end{important}

\section{La posesión como forma de publicidad}

La publicidad por excelencia es la tradición, es decir, la posesión. Rige la máxima \textbf{``sin posesión no hay derecho real''}: la posesión es constitutiva en todos los derechos reales que se ejercen por la posesión. Esto es así porque el derecho real en materia inmobiliaria se adquiere, se transmite y se extingue \textbf{fuera del registro}.

\begin{formula}{Artículo 2006 -- Código Civil y Comercial}
El derecho real inmobiliario se constituye, transmite y extingue por actos jurídicos conforme a lo dispuesto en este Código. La posesión es la forma de ejercer el derecho real.
\end{formula}
% TODO: contenido incompleto o ambiguo en las notas originales (texto del artículo 2006 transcripto tal como figura en el apunte)

\section{Efecto de la inscripción registral}

La inscripción registral es necesaria para la oponibilidad a terceros. Su efecto es \textbf{declarativo}, porque el derecho se constituye, se transmite y se extingue fuera del registro.

Sin inscripción, el derecho será oponible solo a muy pocas personas: aquellas que tuvieron conocimiento del acto (la parte, el escribano, el abogado), ninguna de las cuales podrá prevalerse de la falta de inscripción.

Para la oponibilidad a terceros, en cambio, es imprescindible que el derecho real haya sido inscripto; es decir, que el acto jurídico de adquisición, transmisión, constitución o gravamen del derecho se haya inscripto en el registro. De allí la importancia de la inscripción registral.

% ---------------------------------------------------------------
\chapter{Sistema de publicidad inmobiliaria: Ley 17.801}
\markboth{Sistema de publicidad (Ley 17.801)}{Sistema de publicidad (Ley 17.801)}

Todo lo relativo a la inscripción registral conforma un sistema: el sistema de publicidad inmobiliaria, regulado en Argentina por la \textbf{Ley 17.801}. Se trata de una ley federal y, por lo tanto, aplicable en toda la Nación, sin perjuicio de que los ordenamientos provinciales tengan sus propias normas, en general decretos-leyes, para regular sus propios registros.

\section{Características del sistema registral argentino}

\subsection{Sistema real}

El sistema argentino es un \textbf{sistema real}, porque el registro se realiza \textbf{por inmueble} y no por la persona de su titular. Cada inmueble está matriculado y tiene su correspondiente número de orden. Se le asigna una hoja denominada \textbf{folio real}, donde se asienta todo lo relativo a la adquisición, la transmisión, el gravamen y el levantamiento de gravámenes; es decir, todo aquello que tenga que ver con el derecho real respecto de ese inmueble.

\begin{important}
\textbf{Sistema real vs. sistema personal.} El sistema real (como el argentino) registra por inmueble; el sistema personal registra por titular. Esto permite un seguimiento claro del estado del inmueble, y no de la persona.
\end{important}

\subsection{Sistema de inscripción}

Es un \textbf{sistema de inscripción}, porque cada vez que se inscribe un acto jurídico de adquisición o transmisión del derecho, lo que se asienta se hace en forma sumaria: es una síntesis del documento. En cambio, en los sistemas de transcripción se transcribe el documento completo.

\subsection{Sistema no convalidante}

Es un sistema \textbf{no convalidante}, porque la inscripción no subsana ni convalida el título. Si el título inscripto tiene defectos que pudieran acarrear su nulidad, seguirá teniéndolos aun estando inscripto y será atacable. Otros sistemas, en cambio, validan los títulos y por ello prevén otro procedimiento para atacarlos.

% ---------------------------------------------------------------
\chapter{Principios registrales fundamentales}
\markboth{Principios registrales}{Principios registrales}

Poco tiempo después de la sanción de la Ley 17.801 se realizó un Congreso Internacional de Derecho Registral, en el que se aprobó la \textbf{Carta de Buenos Aires}. Allí se declaran los principios registrales que constituyen la base sobre la que se asienta el sistema registral. Estos principios se explican en su esencia, sin pretender ser una definición concreta de cada principio.

\begin{table}[H]
\centering
\begin{tabularx}{\textwidth}{>{\raggedright\arraybackslash}p{0.27\textwidth}>{\raggedright\arraybackslash}X}
\toprule
\textbf{Principio} & \textbf{Explicación} \\
\midrule
\textbf{Rogación} & El registro procede a solicitud de parte, no de oficio. \\
\textbf{Inscripción} & Documentos inscribibles para su oponibilidad a terceros. \\
\textbf{Presunción registral} & Los asientos registrales se presumen veraces (fe pública). \\
\textbf{Especialidad} & Identificación perfecta e indubitativa del inmueble. \\
\textbf{Legalidad} & Examen de formas extrínsecas (calificación registral). \\
\textbf{Prioridad} & La inscripción primera tiene prioridad sobre las posteriores. \\
\textbf{Tracto sucesivo} & Encadenamiento perfecto entre todos los títulos. \\
\textbf{Publicidad} & Dar publicidad y seguridad al tráfico inmobiliario. \\
\bottomrule
\end{tabularx}
\caption{Principios registrales.}
\end{table}

\section{Principio de rogación}

El \textbf{principio de rogación} significa que el registro procede a solicitud de parte y no de oficio. La intervención del registro es siempre rogada (solicitada), ya sea para lograr un asiento inicial, para las mutaciones registrales, para anotar embargos y anotaciones de litis, o para pedir certificados o informes.

\section{Principio de inscripción}

El \textbf{principio de inscripción} establece que todo documento por el que se constituyen, transmiten, declaran, modifican o extinguen derechos reales sobre inmuebles debe ser inscripto para su oponibilidad. Así lo prevé el artículo 2 de la Ley 17.801, y se vincula con lo ya explicado: el efecto declarativo de la inscripción y su efecto de oponibilidad a terceros.

\begin{formula}{Artículo 2 -- Ley 17.801}
Todo documento por el que se constituyen, transmiten, declaran, modifican o extinguen derechos reales sobre inmuebles debe ser inscripto para su oponibilidad a terceros.
\end{formula}

\section{Principio de presunción registral}

El \textbf{principio de presunción registral} se infiere del artículo 22 de la ley e implica que los asientos registrales se presumen veraces. Se enlaza con la \textbf{fe pública registral}.

Si se obtiene un informe expedido por el registro que indica que un inmueble pertenece a una persona determinada, esa información es el único canal para conocer el derecho real que resulta oponible. Aunque pueda haber existido algún caso de acciones judiciales contra el registro, por eso lo que indica el registro se presume veraz.

\section{Principio de especialidad}

El \textbf{principio de especialidad} es de gran importancia. En todo título que se inscribe, plasmado en una escritura pública, el inmueble debe estar perfectamente individualizado: las medidas, la superficie, los linderos, la nomenclatura, y todo dato que indique indubitablemente respecto de qué inmueble se realiza la inscripción. Esto surge de todo el capítulo tercero de la ley, referido a la matriculación.

\section{Principio de legalidad o calificación}

El \textbf{principio de legalidad}, también llamado \textbf{principio de calificación}, confiere al registro la facultad de examinar las \textbf{formas extrínsecas} del título cuya inscripción se pretende; no su contenido intrínseco. Si las formas extrínsecas reúnen todos los requisitos exigibles, el registro procede a inscribirlo; si no, lo rechaza. Puede rechazarlo rotundamente o inscribirlo \textbf{provisionalmente por un plazo de 180 días}.

En ese plazo debe subsanarse la observación, tanto si la inscripción fue solicitada por el escribano para una escritura pública como si fue requerida por el juzgado mediante la orden de una medida cautelar o de una adjudicación.

\begin{formula}{Artículos 8 y 9 -- Ley 17.801}
Los artículos 8 y 9 contemplan los recursos que pueden interponerse, administrativamente y, agotada la vía administrativa, judicialmente, si se considera que el registro se ha opuesto a la inscripción de manera infundada.
\end{formula}

\section{Principio de prioridad}

El \textbf{principio de prioridad} indica que tiene prioridad el acto que ingresó primero al registro sobre otro acto que ingresó posteriormente, aun cuando el título posterior sea de fecha anterior. Lo que otorga la prioridad es la \textbf{fecha de inscripción y no la fecha del título}.

\section{Principio de tracto sucesivo}

El \textbf{principio de tracto sucesivo} implica que cada transmisión sucede a la otra mediante un \textbf{encadenamiento perfecto entre todos los títulos}: los unos derivan de los otros en forma regular, ordenada y encadenada.

Las \textbf{excepciones} a este principio se dan cuando no existe transmisión de un derecho a otro; por ejemplo, en las inscripciones originarias y en los inmuebles adquiridos por prescripción adquisitiva, que no responden al sistema de tradición traslativa de dominio.

\begin{formula}{Tracto abreviado -- Artículo 16, Ley 17.801}
El tracto abreviado permite que, en algunos casos regulados por la ley, puedan inscribirse simultáneamente dos transmisiones en un mismo asiento. No es una excepción al tracto sucesivo, sino una facilitación que mantiene el encadenamiento registral.
\end{formula}

\begin{example}{Tracto abreviado en la sucesión}
Juan compra un departamento y luego fallece. Sus hijos heredan el departamento y obtienen la declaratoria de herederos. El juzgado ordena inscribir el inmueble a nombre de los herederos. Si los herederos deciden vender el inmueble, resulta costoso realizar dos trámites: primero inscribir a nombre de los herederos y luego inscribir la venta. Aquí aparece el beneficio del artículo 16 de la Ley 17.801: se solicita al juez del sucesorio que la inscripción a nombre de los herederos se realice directamente con la escritura traslativa de dominio por la venta. Es un procedimiento sumamente utilizado.
\end{example}

\section{Principio de publicidad}

El \textbf{principio de publicidad} surge de todo lo expuesto. El objeto fundamental del registro y su finalidad específica es dar publicidad y conferir seguridad al tráfico inmobiliario, y abastece la buena fe de los derechos reales. La buena fe es registral: no puede alegarse buena fe si no se obtuvieron los informes que permiten conocer verazmente la información que consta en el registro.

% ---------------------------------------------------------------
\chapter{Documentos que se inscriben}
\markboth{Documentos inscribibles}{Documentos inscribibles}

Lo que se inscribe siempre son \textbf{documentos}. Aunque en el lenguaje cotidiano se diga que ``se inscribe la transmisión del derecho'', ello siempre está plasmado en un instrumento público, porque se trata de materia inmobiliaria.

\begin{formula}{Artículo 2, Ley 17.801 -- Documentos inscribibles}
Todos aquellos instrumentos que constituyen, transmiten o extinguen derechos reales, también los que dispongan embargos, inhibiciones y las otras providencias cautelares. Debe tratarse de instrumentos públicos, emanados del notario que realiza la escritura pública.
\end{formula}

\section{Instrumentos notariales}

El notario es el profesional cuya incumbencia es la realización de la escritura pública, como fedatario, y tiene a su cargo la inscripción de la transmisión y adquisición del derecho real.

\section{Instrumentos judiciales}

También se inscriben los instrumentos judiciales: las resoluciones judiciales que ordenan o disponen la inscripción. Por ejemplo:

\begin{itemize}
    \item \textbf{Sucesión por causa de muerte:} el inmueble que pertenecía a una persona pasa a ser de los herederos.
    \item \textbf{Liquidación de la sociedad conyugal:} cuando un inmueble se adjudica a uno de los cónyuges, el juez ordena la adjudicación y emite el documento que dispone la inscripción a nombre del adjudicatario.
\end{itemize}

% ---------------------------------------------------------------
\chapter{Bloqueo registral y reserva de prioridad}
\markboth{Bloqueo registral}{Bloqueo registral}

\section{Significado del bloqueo registral}

El \textbf{bloqueo registral}, también llamado \textbf{reserva de prioridad}, es un procedimiento por el cual el registro \textbf{congela la información}.

Todos los escribanos pueden pedir la información a través del \textbf{informe de dominio}, que da cuenta del titular, del porcentaje de titularidad, de los gravámenes y de las restricciones. Sin embargo, cuando un escribano solicita el certificado en función de la realización de un negocio inmobiliario, ese certificado tiene un efecto diferente, una virtualidad que no tiene el informe.

Ese certificado produce la \textbf{prioridad} del negocio jurídico para el cual el escribano lo solicitó. Garantiza la inmutabilidad de todos los datos que el registro da a conocer durante un plazo de \textbf{15 días en la ciudad, 25 días en la provincia o 30 días fuera de la provincia}.

\begin{important}
El informe de dominio solo brinda información. El certificado solicitado para un negocio inmobiliario, en cambio, produce el bloqueo registral: los datos del inmueble se mantienen inmutables durante el plazo de vigencia.
\end{important}

\section{La metáfora de la sombrilla (paraguas)}

El funcionamiento del bloqueo registral se explica mediante una línea de tiempo y una metáfora.

\begin{note}
\textbf{Metáfora de la sombrilla registral.} Durante los 15 días (o 25--30, según jurisdicción) que dura la validez del certificado, se abre una sombrilla protectora sobre el inmueble. Mientras está abierta, los nuevos gravámenes (como gotitas de lluvia) no pueden ingresar. Si la sombrilla se cierra antes de que el escribano haya inscripto la venta, los gravámenes caen sobre el nuevo propietario.
\end{note}

\subsection{El caso planteado}

\begin{example}{Compraventa entre Juan y María}
Juan, el enajenante, adquirió el inmueble por 50 mil dólares. Pacta con María, en un boleto (que es un derecho personal), que ella pagará 20 mil dólares al boleto y 30 mil dólares a la escritura. El boleto se celebra el 10 de junio y se acuerda escriturar el 10 de julio, a 30 días. Se lleva el título de propiedad original al escribano, quien solicita los certificados cuya expedición produce el bloqueo registral.
\end{example}

\begin{table}[H]
\centering
\begin{tabularx}{\textwidth}{>{\raggedright\arraybackslash}p{0.27\textwidth}>{\raggedright\arraybackslash}X}
\toprule
\textbf{Momento} & \textbf{Hecho} \\
\midrule
10 de junio & Celebración del boleto (derecho personal de María). \\
1.º al 15 de julio & Plazo de 15 días de vigencia de los certificados (``paraguas abierto''). \\
10 de julio & Otorgamiento de la escritura traslativa de dominio, dentro del período de vigencia de los certificados. \\
10 de julio al 15 de agosto & Plazo de 45 días para inscribir el título. \\
\bottomrule
\end{tabularx}
\caption{Línea de tiempo del caso planteado.}
\end{table}
% TODO: contenido incompleto o ambiguo en las notas originales (las fechas del apunte no coinciden exactamente con los plazos de 15 y 45 días indicados)

\subsection{Efecto de la expedición del certificado}

La expedición del certificado abre un ``paraguas'' que protege el inmueble: durante ese tiempo, los datos que indican que el inmueble es de Juan, que no está gravado y que no tiene hipotecas se mantienen inalterables.

Si en ese lapso el Juzgado Civil 5 ordena un embargo de 100 mil dólares sobre el inmueble matrícula 3267, titularidad de Juan, el gravamen \textbf{no puede ingresar} mientras el paraguas está abierto: como una gotita, el embargo queda suspendido.

Mientras tanto, es obligación del escribano otorgar la escritura traslativa de dominio, con el título de propiedad a la vista y los certificados vigentes. Hasta ese momento había solo un derecho personal de María; con la escritura pública celebrada dentro del período de vigencia (el 10 de julio) se cumple la tradición traslativa de dominio, pero \textbf{falta la inscripción}.

\subsection{Plazo de 45 días y retroactividad}

Aparece entonces un nuevo plazo de \textbf{45 días}, que la ley concede al escribano para inscribir el título. Si lo inscribe dentro de ese plazo, se produce la \textbf{retroactividad}: la inscripción tiene efecto retroactivo al día de la celebración de la escritura, el 10 de julio.

Cuando el paraguas se cierra, el embargo ordenado por el Juzgado Civil 5 cae como la gotita; pero como la inscripción tiene efecto retroactivo al día de la celebración, el inmueble ya no es de Juan sino de María. Por lo tanto, el embargo es rechazable.

\subsection{Inscripción fuera del plazo de 45 días}

Si el escribano no inscribe dentro de los 45 días (por ejemplo, inscribe el 20 de agosto en lugar de antes del 15 de agosto), pierde el privilegio de la retroactividad que otorga la ley. Es posible inscribir después de los 45 días, pero se pierde la retroactividad.

\begin{important}
Pueden inscribirse los títulos después de los 45 días, pero se pierde la retroactividad. En cambio, no puede otorgarse la escritura si no es dentro de la vigencia de los certificados (dentro de la protección del paraguas).
\end{important}

En el caso planteado, la transmisión de Juan a María es oponible a partir del 20 de agosto. Cuando el paraguas se cerró, el embargo ingresó: María compra el inmueble con un embargo de 100 mil dólares. Se trata de un departamento de 50 mil dólares gravado por 100 mil, de modo que \textbf{lo ha perdido todo}.

Si bien puede iniciarse un juicio contra el escribano y obtenerse una indemnización de daños, esta es escasa, y es probable que la deuda ni siquiera fuera de María. El verdadero desafío es \textbf{prevenir el daño}: repararlo ya tiene sabor a poco. Lo que María quería era una casa; lo necesario es estar presente para ayudar a las personas a seguir adelante.

% ---------------------------------------------------------------
\chapter{Importancia del sistema registral y caducidades}
\markboth{Importancia y caducidades}{Importancia y caducidades}

\section{Importancia del sistema registral en la práctica profesional}

El derecho real es de orden público, y el sistema registral está comprometido con la defensa de ese orden público. Es lo que se debe procurar y prevenir mediante el buen funcionamiento del sistema.

El sistema registral tiene una importancia fundamental en esta rama del derecho porque permite tomar conocimiento de todo lo inscripto: al comprar un inmueble, al embargarlo, o cuando se es acreedor quirografario que ganó el juicio y quiere conocer si el demandado condenado tiene inmuebles a su nombre. Además, brinda \textbf{seguridad al tráfico inmobiliario}.

\section{Caducidades y plazos registrales}

Existen ciertos plazos de caducidad vinculados con la inscripción.

\begin{table}[H]
\centering
\begin{tabularx}{\textwidth}{>{\raggedright\arraybackslash}p{0.38\textwidth}>{\raggedright\arraybackslash}X}
\toprule
\textbf{Inscripción} & \textbf{Caducidad} \\
\midrule
Inscripción provisional & A los 180 días, si no se subsana la observación. \\
Medidas cautelares & A los 5 años. \\
Inscripciones hipotecarias & A los 35 años. \\
\bottomrule
\end{tabularx}
\caption{Caducidades registrales.}
\end{table}

% ---------------------------------------------------------------
\chapter{Resumen Cornell}
\markboth{Resumen Cornell}{Resumen Cornell}

\begingroup
\RaggedRight
\begin{longtable}{>{\RaggedRight\arraybackslash}p{0.27\textwidth}>{\RaggedRight\arraybackslash}p{0.63\textwidth}}
\toprule
\textbf{Preguntas / palabras clave} & \textbf{Notas / respuestas} \\
\midrule
\endfirsthead
\toprule
\textbf{Preguntas / palabras clave} & \textbf{Notas / respuestas} \\
\midrule
\endhead
\midrule
\multicolumn{2}{r}{\small\emph{Continúa en la página siguiente}} \\
\endfoot
\bottomrule
\endlastfoot

\textbf{¿Cuál es la diferencia fundamental entre un derecho real y un derecho personal?} &
El derecho real es absoluto y oponible \emph{erga omnes} (a todos), y debe ser respetado por todos. El derecho personal es relativo y solo es oponible a las partes que conforman el vínculo obligacional. Por ello, el derecho real requiere publicidad (debe ser conocido por todos), mientras que el personal no. \\
\addlinespace
\cmidrule(lr){1-2}

\textbf{¿Cuál es el efecto exacto de la inscripción registral?} &
La inscripción tiene efecto declarativo. El derecho real se constituye, transmite y extingue fuera del registro. La inscripción no lo crea, sino que lo hace oponible a terceros. Sin inscripción, el derecho es oponible solo a quienes tuvieron conocimiento del acto (partes, escribano, abogado). \\
\addlinespace
\cmidrule(lr){1-2}

\textbf{¿Por qué el sistema registral argentino se llama ``sistema real''?} &
Porque el registro se organiza por inmueble y no por persona. Cada inmueble tiene su propia matrícula y folio real, donde se asientan todas sus mutaciones. El sistema personal se organizaría por el titular del inmueble. \\
\addlinespace
\cmidrule(lr){1-2}

\textbf{¿Por qué el sistema es de inscripción y no convalidante?} &
Es de inscripción porque se asienta una síntesis del documento (en los sistemas de transcripción se transcribe el documento completo). Es no convalidante porque la inscripción no subsana ni convalida el título: si tiene defectos que pudieran acarrear su nulidad, seguirá teniéndolos y será atacable. \\
\addlinespace
\cmidrule(lr){1-2}

\textbf{¿Qué establece el artículo 2 de la Ley 17.801?} &
Que todo documento por el que se constituyan, transmitan, declaren, modifiquen o extingan derechos reales sobre inmuebles debe ser inscripto para su oponibilidad a terceros. \\
\addlinespace
\cmidrule(lr){1-2}

\textbf{¿Cuáles son los 8 principios registrales fundamentales?} &
(1) Rogación (procede a solicitud de parte); (2) Inscripción (requisito para la oponibilidad); (3) Presunción registral (los asientos se presumen veraces); (4) Especialidad (identificación indubitativa del inmueble); (5) Legalidad o calificación (examen de formas extrínsecas); (6) Prioridad (la primera inscripción tiene prioridad); (7) Tracto sucesivo (encadenamiento perfecto de títulos); (8) Publicidad (dar seguridad al tráfico inmobiliario). \\
\addlinespace
\cmidrule(lr){1-2}

\textbf{¿Qué examina el registro en virtud del principio de legalidad y qué puede resolver?} &
Examina las formas extrínsecas del título, no su contenido intrínseco. Si reúnen los requisitos, inscribe; si no, rechaza o inscribe provisionalmente por 180 días para subsanar la observación. Los artículos 8 y 9 prevén los recursos administrativos y judiciales. \\
\addlinespace
\cmidrule(lr){1-2}

\textbf{¿Qué otorga la prioridad: la fecha del título o la fecha de inscripción?} &
La fecha de inscripción. Tiene prioridad el acto que ingresó primero al registro, aun cuando el título posterior sea de fecha anterior. \\
\addlinespace
\cmidrule(lr){1-2}

\textbf{¿Qué es el ``tracto abreviado'' y qué permite?} &
El tracto abreviado (artículo 16, Ley 17.801) permite inscribir simultáneamente dos transmisiones en un mismo asiento. Ejemplo: cuando los herederos reciben un inmueble por sucesión y luego lo venden, se pueden inscribir ambas transmisiones juntas, evitando un trámite costoso. Mantiene el encadenamiento registral. \\
\addlinespace
\cmidrule(lr){1-2}

\textbf{¿Cómo funciona el bloqueo registral o reserva de prioridad?} &
Cuando un escribano solicita certificados para una transacción, estos congelan la información del inmueble por 15 días (ciudad), 25 días (provincia) o 30 días (fuera de la provincia). Durante ese período, los nuevos gravámenes (como embargos) no pueden inscribirse. Si el escribano otorga la escritura dentro de ese plazo, tendrá 45 días más para inscribir, con efecto retroactivo. \\
\addlinespace
\cmidrule(lr){1-2}

\textbf{¿Qué sucede si el escribano no inscribe dentro de los 45 días?} &
Pierde la retroactividad. La inscripción será oponible solo desde la fecha de inscripción real y no desde la fecha de celebración del acto. Cualquier embargo o gravamen que haya ingresado en el ínterin será válido y afectará al nuevo propietario. \\
\addlinespace
\cmidrule(lr){1-2}

\textbf{¿Por qué es crítico prevenir daños en transacciones inmobiliarias?} &
Porque reparar el daño posteriormente es insuficiente. El caso de María, que compra un departamento de 50 mil dólares y lo recibe con un embargo de 100 mil dólares, ilustra que perder todo es irreparable. El desafío es estar presente para prevenir, asesorando correctamente y cuidando todos los plazos registrales. \\
\addlinespace
\cmidrule(lr){1-2}

\textbf{¿Cuáles son los plazos de caducidad en las inscripciones registrales?} &
Inscripción provisional: caduca a los 180 días si no se subsana la observación. Medidas cautelares: caducan a los 5 años. Inscripciones hipotecarias: caducan a los 35 años. \\
\midrule

\multicolumn{2}{p{0.92\textwidth}}{%
\textbf{Resumen:} Los derechos reales inmobiliarios son absolutos y oponibles \emph{erga omnes}, lo que exige publicidad: nadie puede ser obligado a respetar lo que no conoce. La Ley 17.801 establece un sistema registral real, de inscripción y no convalidante, que descansa en 8 principios y cuya inscripción tiene efecto declarativo. El bloqueo registral protege el inmueble durante la vigencia del certificado y, si se inscribe dentro de los 45 días, la inscripción es retroactiva; fuera de ese plazo se pierde la retroactividad.} \\
\end{longtable}
\endgroup

\section*{Síntesis integradora}
\addcontentsline{toc}{section}{Síntesis integradora}

Los derechos reales inmobiliarios son derechos absolutos y oponibles \emph{erga omnes}, lo que significa que deben ser respetados por todos. Esta característica genera una necesidad intrínseca de publicidad: no se puede exigir a nadie que respete aquello que no conoce. Por eso, desde la posesión (forma primaria de publicidad) hasta la inscripción registral, el derecho real requiere ser conocido públicamente.

En Argentina, la Ley 17.801 establece un sistema de publicidad inmobiliaria que es \textbf{real} (por inmueble, no por persona), \textbf{de inscripción} (sumaria, no transcripción completa) y \textbf{no convalidante} (no subsana defectos del título). Este sistema descansa en 8 principios registrales: rogación, inscripción, presunción registral, especialidad, legalidad, prioridad, tracto sucesivo y publicidad, cada uno con una función específica en la seguridad del tráfico inmobiliario. El efecto de la inscripción es \textbf{declarativo}, no constitutivo: el derecho existe sin registro, pero solo es oponible a terceros una vez inscripto.

Una herramienta crucial del sistema es el \textbf{bloqueo registral o reserva de prioridad}, que protege los datos del inmueble por 15 a 30 días cuando se solicita un certificado para una transacción. Ofrece el tiempo necesario para otorgar la escritura y luego inscribir con efecto retroactivo, evitando que gravámenes superpuestos afecten al nuevo propietario. Si se inscribe fuera del plazo de 45 días, se pierde esta protección.

El profesional del derecho es responsable de vigilar estos plazos y mecanismos, previniendo daños que después serán irreparables. El sistema no es perfecto, pero es eficaz cuando se lo usa correctamente.

\end{document}
